% Written by Anders Sjoqvist and Ulf Lundstrom, 2009
% The main sources are: tinyKACTL, Beta and Wikipedia

\chapter{Toán học}

\section{Phương trình}
\[ax^2+bx+c=0 \Rightarrow x = \frac{-b\pm\sqrt{b^2-4ac}}{2a}\]

Cực trị đạt tại $x = -b/2a$.

\[\begin{aligned}ax+by=e\\cx+dy=f\end{aligned}
\Rightarrow
\begin{aligned}x=\dfrac{ed-bf}{ad-bc}\\y=\dfrac{af-ec}{ad-bc}\end{aligned}\]

Tổng quát, với hệ phương trình $Ax = b$ (định lý Cramer), nghiệm của biến $x_i$ là:
\[x_i = \frac{\det A_i'}{\det A} \]
trong đó $A_i'$ là ma trận $A$ với cột thứ $i$ được thay thế bởi vector $b$.

\section{Hệ thức truy hồi}
Nếu $a_n = c_1 a_{n-1} + \dots + c_k a_{n-k}$, và $r_1, \dots, r_k$ là các nghiệm phân biệt của đa thức đặc trưng $x^k - c_1 x^{k-1} - \dots - c_k = 0$, thì tồn tại $d_1, \dots, d_k$ sao cho:
\[a_n = d_1r_1^n + \dots + d_kr_k^n. \]
Nghiệm bội $r$ sẽ có nhân tử đa thức, ví dụ $a_n = (d_1n + d_2)r^n$.

\section{Lượng giác}
\begin{align*}
\sin(v+w)&{}=\sin v\cos w+\cos v\sin w\\
\cos(v+w)&{}=\cos v\cos w-\sin v\sin w\\
\end{align*}
\begin{align*}
\tan(v+w)&{}=\dfrac{\tan v+\tan w}{1-\tan v\tan w}\\
\sin v+\sin w&{}=2\sin\dfrac{v+w}{2}\cos\dfrac{v-w}{2}\\
\cos v+\cos w&{}=2\cos\dfrac{v+w}{2}\cos\dfrac{v-w}{2}
\end{align*}
\[ (V+W)\tan(v-w)/2{}=(V-W)\tan(v+w)/2 \]
trong đó $V, W$ là độ dài các cạnh đối diện góc $v, w$.
\begin{align*}
	a\cos x+b\sin x&=r\cos(x-\phi)\\
	a\sin x+b\cos x&=r\sin(x+\phi)
\end{align*}
với $r=\sqrt{a^2+b^2}, \phi=\operatorname{atan2}(b,a)$.

\section{Hình học}

\subsection{Tam giác}
Độ dài các cạnh: $a,b,c$\\
Nửa chu vi: $p=\dfrac{a+b+c}{2}$\\
Diện tích (Heron): $A=\sqrt{p(p-a)(p-b)(p-c)}$\\
Bán kính đường tròn ngoại tiếp: $R=\dfrac{abc}{4A}$\\
Bán kính đường tròn nội tiếp: $r=\dfrac{A}{p}$\\
Độ dài trung tuyến: $m_a=\tfrac{1}{2}\sqrt{2b^2+2c^2-a^2}$\\
Độ dài phân giác: $s_a=\sqrt{bc\left[1-\left(\dfrac{a}{b+c}\right)^2\right]}$\\
Định lý sin: $\dfrac{\sin\alpha}{a}=\dfrac{\sin\beta}{b}=\dfrac{\sin\gamma}{c}=\dfrac{1}{2R}$\\
Định lý cosin: $a^2=b^2+c^2-2bc\cos\alpha$\\
Định lý tang: $\dfrac{a+b}{a-b}=\dfrac{\tan\dfrac{\alpha+\beta}{2}}{\tan\dfrac{\alpha-\beta}{2}}$\\

\subsection{Tứ giác}
Với độ dài các cạnh $a,b,c,d$, hai đường chéo $e, f$, góc giữa hai đường chéo $\theta$, diện tích $A$ và
$F=b^2+d^2-a^2-c^2$:

\[ 4A = 2ef \cdot \sin\theta = F\tan\theta = \sqrt{4e^2f^2-F^2} \]

Với tứ giác nội tiếp, tổng hai góc đối bằng $180^\circ$, định lý Ptolemy $ef = ac + bd$, và công thức Brahmagupta $A = \sqrt{(p-a)(p-b)(p-c)(p-d)}$.

\subsection{Tọa độ cầu}
\begin{center}
\includegraphics[width=25mm]{content/math/sphericalCoordinates}
\end{center}
\[\begin{array}{cc}
x = r\sin\theta\cos\phi & r = \sqrt{x^2+y^2+z^2}\\
y = r\sin\theta\sin\phi & \theta = \textrm{acos}(z/\sqrt{x^2+y^2+z^2})\\
z = r\cos\theta & \phi = \textrm{atan2}(y,x)
\end{array}\]

\section{Đạo hàm / Tích phân}
\begin{align*}
	\dfrac{d}{dx}\arcsin x = \dfrac{1}{\sqrt{1-x^2}} &&& \dfrac{d}{dx}\arccos x = -\dfrac{1}{\sqrt{1-x^2}} \\
	\dfrac{d}{dx}\tan x = 1+\tan^2 x &&& \dfrac{d}{dx}\arctan x = \dfrac{1}{1+x^2} \\
	\int\tan ax = -\dfrac{\ln|\cos ax|}{a} &&& \int x\sin ax = \dfrac{\sin ax-ax \cos ax}{a^2} \\
	\int e^{-x^2} = \frac{\sqrt \pi}{2} \text{erf}(x) &&& \int xe^{ax}dx = \frac{e^{ax}}{a^2}(ax-1)
\end{align*}

Tích phân từng phần:
\[\int_a^bf(x)g(x)dx = [F(x)g(x)]_a^b-\int_a^bF(x)g'(x)dx\]

\section{Tổng}
\[ c^a + c^{a+1} + \dots + c^{b} = \frac{c^{b+1} - c^a}{c-1}, c \neq 1 \]
\begin{align*}
	1 + 2 + 3 + \dots + n &= \frac{n(n+1)}{2} \\
	1^2 + 2^2 + 3^2 + \dots + n^2 &= \frac{n(2n+1)(n+1)}{6} \\
	1^3 + 2^3 + 3^3 + \dots + n^3 &= \frac{n^2(n+1)^2}{4} \\
	1^4 + 2^4 + 3^4 + \dots + n^4 &= \frac{n(n+1)(2n+1)(3n^2 + 3n - 1)}{30} \\
\end{align*}

\section{Chuỗi khai triển}
$$e^x = 1+x+\frac{x^2}{2!}+\frac{x^3}{3!}+\dots,\,(-\infty<x<\infty)$$
$$\ln(1+x) = x-\frac{x^2}{2}+\frac{x^3}{3}-\frac{x^4}{4}+\dots,\,(-1<x\leq1)$$
$$\sqrt{1+x} = 1+\frac{x}{2}-\frac{x^2}{8}+\frac{2x^3}{32}-\frac{5x^4}{128}+\dots,\,(-1\leq x\leq1)$$
$$\sin x = x-\frac{x^3}{3!}+\frac{x^5}{5!}-\frac{x^7}{7!}+\dots,\,(-\infty<x<\infty)$$
$$\cos x = 1-\frac{x^2}{2!}+\frac{x^4}{4!}-\frac{x^6}{6!}+\dots,\,(-\infty<x<\infty)$$

\section{Lý thuyết xác suất}
Gọi $X$ là biến ngẫu nhiên rời rạc với xác suất nhận giá trị $x$ là $p_X(x)$. Kỳ vọng (mean) $\mu=\mathbb{E}(X)=\sum_xxp_X(x)$ và phương sai $\sigma^2=V(X)=\mathbb{E}(X^2)-(\mathbb{E}(X))^2=\sum_x(x-\mathbb{E}(X))^2p_X(x)$ với $\sigma$ là độ lệch chuẩn. Nếu $X$ liên tục với hàm mật độ xác suất $f_X(x)$, các tổng trên chuyển thành tích phân tương ứng.

Tính tuyến tính của kỳ vọng:
\[\mathbb{E}(aX+bY) = a\mathbb{E}(X)+b\mathbb{E}(Y)\]
Với $X$ và $Y$ độc lập: \[V(aX+bY) = a^2V(X)+b^2V(Y).\]

\subsection{Phân phối rời rạc}

\subsubsection{Phân phối nhị thức}
Số lần thành công trong $n$ phép thử độc lập có xác suất thành công $p$ là $\textrm{Bin}(n,p),\,n=1,2,\dots,\, 0\leq p\leq1$.
\[p(k)=\binom{n}{k}p^k(1-p)^{n-k}\]
\[\mu = np,\,\sigma^2=np(1-p)\]
$\textrm{Bin}(n,p)$ xấp xỉ $\textrm{Po}(np)$ khi $p$ nhỏ.

\subsubsection{Phân phối hình học (First success)}
Số phép thử cần thiết để đạt thành công đầu tiên với xác suất $p$ là $\textrm{Fs}(p),\,0\leq p\leq1$.
\[p(k)=p(1-p)^{k-1},\,k=1,2,\dots\]
\[\mu = \frac1p,\,\sigma^2=\frac{1-p}{p^2}\]

\subsubsection{Phân phối Poisson}
Số sự kiện xảy ra trong khoảng thời gian $t$ với tốc độ trung bình $\kappa$ và độc lập với thời điểm xảy ra sự kiện trước là $\textrm{Po}(\lambda),\,\lambda=t\kappa$.
\[p(k)=e^{-\lambda}\frac{\lambda^k}{k!}, k=0,1,2,\dots\]
\[\mu=\lambda,\,\sigma^2=\lambda\]

\subsection{Phân phối liên tục}

\subsubsection{Phân phối đều}
Nếu hàm mật độ xác suất là hằng số giữa $a$ và $b$ và bằng 0 ở nơi khác: $\textrm{U}(a,b),\,a<b$.
\[f(x) = \left\{
\begin{array}{cl}
\frac{1}{b-a} & a<x<b\\
0 & \textrm{khác}
\end{array}\right.\]
\[\mu=\frac{a+b}{2},\,\sigma^2=\frac{(b-a)^2}{12}\]

\subsubsection{Phân phối mũ}
Thời gian giữa các sự kiện trong quá trình Poisson: $\textrm{Exp}(\lambda),\,\lambda>0$.
\[f(x) = \left\{
\begin{array}{cl}
\lambda e^{-\lambda x} & x\geq0\\
0 & x<0
\end{array}\right.\]
\[\mu=\frac{1}{\lambda},\,\sigma^2=\frac{1}{\lambda^2}\]

\subsubsection{Phân phối chuẩn}
Đa số giá trị ngẫu nhiên thực có kỳ vọng $\mu$ và phương sai $\sigma^2$ được mô tả tốt bởi $\mathcal{N}(\mu,\sigma^2),\,\sigma>0$.
\[ f(x) = \frac{1}{\sqrt{2\pi\sigma^2}}e^{-\frac{(x-\mu)^2}{2\sigma^2}} \]
Nếu $X_1 \sim \mathcal{N}(\mu_1,\sigma_1^2)$ và $X_2 \sim \mathcal{N}(\mu_2,\sigma_2^2)$ thì:
\[ aX_1 + bX_2 + c \sim \mathcal{N}(\mu_1+\mu_2+c,a^2\sigma_1^2+b^2\sigma_2^2) \]

\section{Chuỗi Markov}
\emph{Chuỗi Markov} là một quá trình ngẫu nhiên rời rạc mà trạng thái tiếp theo chỉ phụ thuộc vào trạng thái hiện tại.
Gọi $X_1,X_2,\ldots$ là dãy các biến ngẫu nhiên sinh bởi quá trình Markov.
Ma trận chuyển trạng thái $\mathbf{P} = (p_{ij})$ với $p_{ij} = \Pr(X_n = i | X_{n-1} = j)$,
và $\mathbf{p}^{(n)} = \mathbf P^n \mathbf p^{(0)}$ là phân phối xác suất cho $X_n$,
trong đó $\mathbf{p}^{(0)}$ là phân phối ban đầu.

Phân phối dừng $\mathbf{\pi}$ thỏa mãn $\mathbf{\pi} = \mathbf{\pi P}$.
Nếu chuỗi Markov \emph{không thể phân rã (irreducible)} (có thể đi từ bất kỳ trạng thái nào đến trạng thái nào),
thì $\pi_i = \frac{1}{\mathbb{E}(T_i)}$ với $\mathbb{E}(T_i)$ là thời gian kỳ vọng giữa hai lần ghé thăm trạng thái $i$.
$\pi_j/\pi_i$ là số lần kỳ vọng ghé thăm trạng thái $j$ giữa hai lần ghé thăm trạng thái $i$.

Với đồ thị vô hướng liên thông không hai phía, xác suất chuyển đều giữa các đỉnh kề, $\pi_i$ tỉ lệ thuận với bậc của đỉnh $i$.

Một chuỗi Markov có tính \emph{ergodic} nếu phân phối tiệm cận độc lập với phân phối ban đầu.
Chuỗi Markov hữu hạn có tính ergodic khi và chỉ khi nó irreducible và \emph{không tuần hoàn (aperiodic)} (ước chung lớn nhất của độ dài các chu trình bằng 1).
$\lim_{k\rightarrow\infty}\mathbf{P}^k = \mathbf{1}\pi$.

Một chuỗi Markov là A-chain nếu các trạng thái có thể chia thành hai tập $\mathbf{A}$ và $\mathbf{G}$, sao cho mọi trạng thái trong $\mathbf{A}$ là trạng thái hấp thụ ($p_{ii}=1$), và mọi trạng thái trong $\mathbf{G}$ đều dẫn tới một trạng thái hấp thụ trong $\mathbf{A}$.
Xác suất hấp thụ tại trạng thái $i\in\mathbf{A}$ khi bắt đầu từ $j$ là $a_{ij} = p_{ij}+\sum_{k\in\mathbf{G}} a_{ik}p_{kj}$.
Thời gian kỳ vọng đến khi bị hấp thụ từ $i$ là $t_i = 1+\sum_{k\in\mathbf{G}}p_{ki}t_k$.
